\chapter{Research Methodology}
\label{chap:methodology-en}

This chapter translates the theoretical constructs of Chapter~\ref{ch:related-work}
into observables, decision rules, and admission boundaries. Every principal claim
preserves the same path: theoretical construct $\rightarrow$ operational definition
$\rightarrow$ experiment $\rightarrow$ result $\rightarrow$ interpretation. The
C01--C11 correspondence is additionally registered in
\code{thesis/data/thesis\_traceability.json}.

\section{Principles of the Research Design}

The methodology is organized around four principles: prior specification, separation
of data roles, preservation of negative results, and reproducible artifact provenance.
They are not imposed as one immutable protocol over the entire multi-month program;
rather, they form a sequence of versioned contracts. Completed Stage~1/2 work is not
rewritten when Stage~3 opens, and later mechanistic-attribution and QWake questions
receive their own freeze and admission boundaries.

The Stage~1/2 specification fixes the primary statistical unit, datasets, principal
contrasts, registered operational equivalence margin, pilot-selection rules, and the
first admissible point of access to the final test set. Stage~3 is added through a
separate protocol. SI-MA0, SI-MA1, B1/B2, and QWake inherit established results as
predecessor evidence but do not redefine them retrospectively.

The experiment sequence is therefore treated as a directed claim chain. A new stage
may support, refine, or restrict interpretation of an earlier result, but it cannot
make an unsuccessful or negative predecessor ``cease to exist''.

\section{Evolution of Research Scope and Historical Branches}
\label{sec:method-scope-evolution-en}

The root \code{HYPOTHESES.md} and \code{ROADMAP.md} preserve the history of a broader
research program and therefore contain historical RQ1--RQ11, Scenario~A,
PC-TREF/PC-CATM, and several exploratory directions. These identifiers are not in
one-to-one correspondence with the final dissertation RQ1--RQ3. The final questions
consolidate only branches that form a completed evidence chain or were explicitly
narrowed by a later protocol.

\begin{longtable}{@{}L{0.19\textwidth}L{0.31\textwidth}L{0.39\textwidth}@{}}
\toprule
Historical branch & Final role & Dissertation status \\
\midrule
\endfirsthead
\toprule
Historical branch & Final role & Dissertation status \\
\midrule
\endhead
RQ1: quality & RQ1, Stage~1/2 & Included: paired macro-F1 equivalence and the cross-version quality/runtime surface. \\
RQ2: gradients; RQ3: representations & RQ1, Stage~3A & Included as layer-wise gradient analysis and CKA/RSA analysis. \\
RQ4: cost; RQ6: locality; RQ7: organization of exact computation & RQ2, Stage~3B & Included through B0, SI-MA0/1, B1/B2, and matched profiling. \\
RQ5: robustness to noise, blur, and occlusion & Separate early branch & Not included in the final completed evidence chain; no independent dissertation claim is made. \\
RQ8: adaptive budget & RQ3, QWake-FP & Narrowed: instead of general adaptive stopping, the work tests task-relative admissibility, pre-action recognizability with zero observed dangerous accepts, and economics under frozen accounting. \\
RQ9: periodic VJP updates & Historical planned family & Not included in the final confirmatory chain; absence of a result is not treated as a negative experiment. \\
RQ10: approximate feedback & Conditional exploratory branch & Outside the final primary claim contract. \\
RQ11: predict--correct acceleration & Design-history/ADR branch & Not formulated as a completed confirmatory dissertation claim. \\
PC-TREF/PC-CATM and QWake-PC & Theoretical framework for RQ3 & Included: PC-TREF defines task-relative equivalence and representation sufficiency; PC-CATM defines mechanistic feature sources; QWake-PC links them to computational action. Only bounded QWake-FP is evaluated empirically. \\
\bottomrule
\end{longtable}

This reconciliation is important when interpreting repository history. A branch absent
from the final claim contract is not automatically a rejected hypothesis; it may be a
historically planned or exploratory direction, replaced by a narrower estimand, or
left outside the completed evidence chain. Only results explicitly bound to final
RQ1--RQ3 and the machine-readable claim contract are used as dissertation conclusions.

\section{Statistical Unit and Repeated Measurements}
\label{sec:method-stat-unit-en}

For Stage~1/2 and the corresponding Stage~3 comparisons, the independent statistical
unit is an independently trained model identified by \texttt{model\_seed}. Batches,
layers, states, tensor elements, and repeated measurements within one model are not
independent replications.

This distinction is essential for layer-wise analysis. Stage~3A creates many layer and
metric observations, but confirmatory comparisons are first aggregated at model level
and only then paired across matching \texttt{model\_seed} values. Linear depth slopes
and cross-layer matrices may be used descriptively but do not increase the number of
independent models.

The original Stage~1/2 protocol requires at least ten complete pairs. A recommendation
procedure could increase, but not decrease, the number before final test access. The
confirmatory Stage~3A uses \texttt{model\_seed}=0--9; this limit is propagated into
the inference boundary.

\section{Datasets, Architecture, and Computational Domain}

FashionMNIST is the primary Stage~1/2 dataset and MNIST the secondary dataset. The
primary architecture is \texttt{lenet\_classic}. Principal contrasts compare
FixedPred and Strict with BP, while Exact serves as a technical and secondary control.
These conditions define the empirical domain of the dissertation and are not claimed
to be universal properties of PC.

Stage~3A retains FashionMNIST, \texttt{lenet\_classic}, and ten model seeds but moves
from final metrics to layer-wise gradients and representations. Stage~3B uses a
separate frozen ROCm/float32 environment and execution matrices designed for
mechanism, runtime, and memory measurements. A change of analytical level therefore
receives a new protocol and freeze rather than silently extending an earlier estimand.

\section{Separation of Pilot, Validation, Test, and Post-action Oracle Roles}
\label{sec:method-data-separation-en}

In Stage~1/2, the pilot phase uses validation data. The final test set becomes
available only after the \texttt{pilot-freeze}, which records commit, computational
environment, selected parameters, configuration hashes, preregistration, and analysis
plan. Configuration changes after inspecting the final test set do not replace the
original confirmatory analysis and require a separate exploratory identifier.

QWake adds a second distinction between \emph{pre-action input} and
\emph{post-action oracle verification}. A rule may use only the preregistered
representation available before the action. The oracle label is produced after the
action to audit sufficiency and evaluate dangerous accepts, but it is not an input
feature of the rule. Recognizability is therefore tested without direct access by the
decision rule to the target label against which action admissibility is subsequently
checked.

\section{Stage~1/2: Comparative Protocol}
\label{sec:method-stage12-en}

Stage~1/2 establishes the initial comparison surface. Each independently trained model
is compared using predefined metrics and contrasts. The primary quality metric is
macro-F1. In equivalence testing, the equivalence margin is ordinarily tied to the
smallest effect meaningful for the research question \cite{lakens2017equivalence}.
Here, an operational margin of \(\delta=0.01\) is preregistered for the paired
difference between a method and BP. It is a protocol tolerance, not a universal
practical-significance threshold for macro-F1.

For model seed \(s\), define the paired difference
\[
  d_s=m_{\mathrm{PC},s}-m_{\mathrm{BP},s}.
\]
Failure to detect a difference from zero is not treated as equivalence. Equivalence is
admissible only with at least ten complete pairs and when the 90\% confidence interval
for the mean paired difference lies entirely within
\[
  [-\delta,\delta]=[-0.01,0.01].
\]

For primary contrasts the work reports model-level values, mean paired difference,
95\% confidence interval, Cohen \(d_z\), an exact sign-permutation test for a small
number of non-zero pairs or the preregistered Wilcoxon approximation, and Holm
correction for primary contrasts. Difference testing and practical-equivalence testing
therefore serve different logical functions.

Failed and repeated runs are retained append-only. A later technical rerun does not
turn the first terminal failure into a successful independent replication and does not
increase \(n\).

\section{Stage~3A: Layer-wise Confirmatory Analysis}
\label{sec:method-stage3a-en}

Stage~3A tests how well final behavior reflects internal computational structure.
Gradient diagnostics include cosine similarity, norm ratio, relative \(L_2\) error,
and sign agreement; representation diagnostics use CKA and RSA. Comparisons are made
against registered BP references. Exact has a separate numerical-control tolerance.

Confirmatory inference is based on model-level aggregates. For depth dependence,
Spearman correlation is computed separately within each trained model and inference
is then performed over the set of model-level coefficients. Cross-layer CKA matrices
are retained descriptively rather than treated as an additional family of independent
confirmatory tests.

The stage therefore does not ask whether algorithms are simply ``the same''. It asks
which gradient and representation components are preserved or changed in the
registered implementation and configuration.

\section{Stage~3B: From Profiling to Mechanistic Attribution}
\label{sec:method-stage3b-en}

Stage~3B moves the primary estimand to computational mechanism and the cost vector. B0
freezes the canonical ROCm/float32 baseline for FixedPred and Strict. Measurements
include time, memory, and structural execution decomposition, while mathematical
locality and execution locality remain distinct properties.

After B0, the question becomes whether a representation of \texttt{state\_inference}
can be informative enough for a decision about subsequent exact sweeps while remaining
computationally economical. This is addressed through a sequential mechanistic-
attribution program.

\subsection{SI-MA0 and SI-MA1}

SI-MA0 tests the registered primary attribution model. Passing some structural or
recognition checks while failing \texttt{COST-MA0} is not reclassified as an overall
positive result: the negative criterion is preserved and motivates a separate
calibration rather than retrospective threshold adjustment.

SI-MA1 independently evaluates observer-cost calibration. A signed residual is
permitted as a diagnostic quantity; a negative residual is interpreted as
\emph{over-closure} relative to the calibration equation, not negative physical cost
and not ``free'' budget for a future controller. This distinction is carried into
B1/B2 and QWake accounting.

\subsection{B1/B2 and Matched Profiling}

After SI-MA1, exact implementation candidates receive separate preregistered
contracts. B1 is \texttt{isolated\_layer\_vjp}; B2 is \texttt{composite\_vjp}. Each
candidate must pass its own equivalence gate before entering the shared comparison
matrix. No prior direction is assumed for the tradeoff among locality, runtime, and
memory.

Matched B0/B1/B2 profiling is executed as 288 cells corresponding to 96 matched
blocks. The purpose is to compare alternative computational organizations on one
registered surface without mixing the equivalence tolerance with the later engineering
comparison.

\subsection{Operationalizing PC-TREF and PC-CATM in QWake-FP}
\label{sec:method-task-relative-sufficiency-en}

The operationalization preserves theoretical role separation. PC-CATM motivates
mechanistic observables; PC-TREF defines a testable sufficiency relation relative to
a computational action; QWake-FP defines a bounded experimental decision mechanism.
Observing a feature is therefore not permission to take the early action, and action
admissibility is not evidence of positive savings.

The theoretical definitions in Sections~\ref{sec:theory-pc-tref}--
\ref{sec:theory-tref-catm-link} are not redefined here. The QWake-FP branch does not
test full sufficiency of the entire diagnostic representation. It tests the narrower
selective question of whether a rule over registered pre-action FixedPred state can
identify a non-zero region on which the registered analytic completion passes the
post-action reference comparison.

In PC-TREF terms, the complete canonical suffix
\nolinkurl{complete_suffix_stage2_baseline_v1} is the exact counterfactual reference,
and the registered early action
\nolinkurl{fixedpred_eta1_wavefront_completion_v1} belongs to the
\texttt{ANALYTIC\_COMPLETION} family. Both branches start from disposable forks of
the same snapshot; required responses are materialized and the registered comparator
checks the analytic candidate against the exact suffix. \texttt{EARLY\_ADMISSIBLE}
means that the candidate passes this comparison, whereas
\texttt{CONTINUE\_REQUIRED} means it does not. Thus the early action replaces the
remaining canonical iterative sweeps with analytic computation; it does not mean zero
further computation. The post-action label defines the required decision class but is
unavailable to the rule before action. The exact suffix also remains the fallback and
does not collapse the two required classes.

Pre-action features form \(\phi_I(x)\). For a rule \(\pi\), its early-accept region
\(A_{\pi}(a_{\mathrm{early}})\) is checked against the admissible-action region
\(S_{a_{\mathrm{early}}}\). In C2, zero observed dangerous accepts refers only to
the 756 frozen calibration records and is not a statistical population-risk guarantee.

PC-CATM motivates a subset of possible features: channel activity and resulting
correction, mutual cancellation, and inter-layer transport can distinguish states that
look identical under a simple residual norm. C2, however, does not test PC-CATM as a
universal model and need not use its complete feature set. The rule family is frozen
in advance and the result is bounded to that family.

QWake-FP cost accounting follows the same semantics. Runtime of the analytic candidate
is included in the registered compute-cost component together with observation and
control components. What is avoided is the exact iterative suffix, not all
post-decision computation. C09 therefore compares replacement of one computational
path by another rather than comparing ``zero-cost stopping'' with full continuation.

Thus, existence of task-admissible preterminal states, selective recognizability of a
region with zero observed dangerous accepts on the calibration surface, and economic
viability remain three sequential tests. Zero observed dangerous accepts checks only
the finite observed manifestation of
\(A_{\pi}(a_{\mathrm{early}})\subseteq S_{a_{\mathrm{early}}}\); full PC-TREF
sufficiency on the registered \(\Omega\) is not inferred.

\section{QWake-FP: C1 Calibration Surface and Bounded C2 Rule Search}
\label{sec:method-qwake-en}

QWake-FP is the experimental specialization of QWake-PC. After Stage~3B shows the
importance of observation cost and \texttt{state\_inference} structure, the problem
is split into three tests: whether the sealed trajectory surface contains preterminal
records for which the registered analytic completion is admissible relative to the
required response of the exact suffix; whether a non-zero subset can be recognized
from already available pre-action state with zero observed dangerous accepts; and
whether that decision pays for itself relative to the residual exact computation.

\subsection{C1: Sealed Trajectory Surface}

C1 creates calibration evidence containing pre-action observations, measured cost
edges, residual-computation cost, and post-action oracle labels. Its scientific role
is not to select a rule but to create an immutable surface on which C2 can search
without changing trajectories after seeing rule-selection outcomes.

The registered C1 time axis includes decision epochs \(t\in[0,K_{\mathrm{ref}}]\).
Let \(\rho_{j,t}\in\mathcal R_{C1}\) denote a record and
\(x(\rho_{j,t})\in\mathcal X\) its contained pre-action state. The terminal boundary
\(t=K_{\mathrm{ref}}\) therefore belongs to the calibration record surface with zero
remaining sweeps. The set of state values
\(\Omega_{C1}:=\{x(\rho):\rho\in\mathcal R_{C1}\}\) is not identified with the set
of records and need not have the same cardinality. Structurally preterminal records
satisfy \(t<K_{\mathrm{ref}}\), equivalently
\texttt{remaining\_sweeps > 0}; define
\[
  \mathcal R_{\mathrm{pre}}
  =\{\rho_{j,t}\in\mathcal R_{C1}:t<K_{\mathrm{ref}}\}.
\]
Registered C2 coverage is defined over the complete record surface
\(\mathcal R_{C1}\), including the terminal boundary, and is not the fraction of
preterminal applications of the registered analytic action. This distinction is used
only to interpret the already frozen surface: terminal records are not removed and the
C2 estimand and cost model are not recomputed.

C1 data provenance and provenance of later executors are distinct. A new executor must
prove compatibility with predecessor evidence but does not rewrite the provenance of
already sealed C1 artifacts.

\subsection{C2: Frozen Rule Family}

C2 evaluates a finite family of 2,625 frozen scalar rules. It includes a complete-
rejection rule and one-dimensional threshold rules over preregistered numerical
features available before action; threshold search is bounded to values frozen before
C2 execution. The post-action oracle, measured outcomes, and confirmatory or test-stage
data are not admissible inputs to rule generation.

For a rule \(\pi\), let
\begin{itemize}
  \item \(n_{\mathrm{acc}}(\pi)\) be the number of accepted records;
  \item \(n_{\mathrm{danger}}(\pi)\) be the number of dangerous accepts---records for which the rule proposes the early action but the post-action oracle assigns \texttt{CONTINUE\_REQUIRED};
  \item \(n_{\mathrm{eval}}\) be the total number of evaluated records;
  \item \(\Delta_{\mathrm{net}}(\pi)\) be aggregate net saving.
\end{itemize}
Coverage is
\[
  \operatorname{cov}(\pi)=\frac{n_{\mathrm{acc}}(\pi)}{n_{\mathrm{eval}}}.
\]

For a sealed record \(i\), the frozen accounting model defines
\[
  g_i(\pi)=
  \begin{cases}
    \tau_i^{\mathrm{rem}}-\kappa_i,&\text{if the early action is accepted and not dangerous},\\
    -\kappa_i,&\text{otherwise},
  \end{cases}
\]
where \(\tau_i^{\mathrm{rem}}\) is registered residual-computation time and
\(\kappa_i\) is full decision cost aggregated from measured edges. Then
\[
  \Delta_{\mathrm{net}}(\pi)=\sum_i g_i(\pi).
\]

The temporal part of \(\kappa_i\) is decomposed into frozen categories
\[
  \kappa_i=\kappa_i^{\mathrm{compute}}+\kappa_i^{\mathrm{latency}}+
  \kappa_i^{\mathrm{diagnostic}}+\kappa_i^{\mathrm{observer}}+
  \kappa_i^{\mathrm{control}}+\kappa_i^{\mathrm{fallback}}.
\]
Memory is a separate component of the cost vector and is not added to nanoseconds.

A rule is eligible for freezing only when
\[
  n_{\mathrm{danger}}(\pi)=0,\qquad
  n_{\mathrm{acc}}(\pi)>0,\qquad
  \Delta_{\mathrm{net}}(\pi)>0,
\]
and its result class is \texttt{SAFE\_AND\_BENEFICIAL}. If this set is empty, C2
returns a bounded negative result and does not create \texttt{selected-policy.json}.

\subsection{Interpretation Boundary of Full Decision-cost Accounting}

C2 extracts decision cost from sealed C1 \texttt{measured\_edges}; the same per-record
cost surface is therefore applied to every rule. This provides a controlled comparison
and prevents post-outcome revaluation. It also fixes the estimand boundary: C2 tests
economic viability under \emph{frozen full decision-cost accounting}, not the
separately measured marginal runtime overhead of a specific simple recognizer.

Consequently, the negative C2 result is not used as evidence that a minimal
implementation of the simple rule is itself economically non-viable. That is a
separate question with status \texttt{not\_tested}.

\section{Negative Results and the C3 Boundary}
\label{sec:method-negative-en}

A negative result is terminal for its frozen protocol and is not an automatic reason
to broaden the search. After C2, no additional threshold search is performed, the cost
model is not changed, and features are not added to obtain a desired positive outcome
within the same campaign.

C3 is a confirmatory stage only for a rule established by a C2 rule-freeze receipt. A
bounded negative C2 result leaves the \texttt{protocol\_receipt\_kind} value
corresponding to a rule freeze absent, so the receipt chain does not open C3. Absence of
C3 is therefore a consequence of preregistered protocol logic, not a missing positive
experiment.

Historical infrastructure failures are likewise not reinterpreted as scientific
results. The one-shot authorization/admission model distinguishes failure before host
claim, terminal failure after authorization consumption, and successfully sealed
scientific execution. Infrastructure repair requires a new provenance line and new
admission rather than silent reuse of consumed authorization.

\section{Artifact Provenance and Reproducibility}
\label{sec:method-provenance-en}

For every significant stage, available source and environment identities are recorded:
Git commit and tree, environment state or container-image digest, runtime manifest,
request hash, authorization binding, host claim, terminal receipt, and artifact
checksums. Where a protocol requires one-shot execution, automatic retry is forbidden
and authorization consumption is retained independently of later scientific
materialization success.

This scheme serves two purposes. First, it binds a scientific result to a particular
executable state. Second, it prevents a repaired later implementation from silently
claiming provenance of earlier evidence. This is essential in a long sequential
program where the scientific predecessor can remain valid even if the executor of a
later stage changes for infrastructure reasons.

\section{Thesis-facing Evidence and CI Build}
\label{sec:method-thesis-ci-en}

The complete evidence set is not copied into LaTeX as manually maintained numbers.
The dissertation uses a small Git-tracked layer. The file
\texttt{thesis/data/research\_claims.json} stores research question, status, scope,
and evidence link for principal claims. For final QWake C2,
\texttt{thesis/data/qwake\_c2\_verified\_summary.json} stores independently verified
aggregates and frozen hashes and is explicitly marked as not creating new scientific
evidence.

Before XeLaTeX, CI checks schemas, arithmetic identities, and protocol boundaries and
then deterministically renders thesis-facing LaTeX assets. The PDF is built from
Git-tracked source in a separate process, receives SHA-256, and is retained as a CI
artifact. This build is read-only relative to scientific evidence: it may not execute
C1/C2, change cost accounting, regenerate rules, or open C3.

Document reproducibility is thus distinct from experimental reproducibility but linked
through verifiable identities. The dissertation can be rebuilt after editorial
changes without new scientific execution, while scientific claim-contract changes
remain visible as ordinary Git differences.

\section{Methodological Limitations}
\label{sec:method-limitations-en}

Ten model seeds do not turn layer-wise observations into additional independent
replications. FashionMNIST and \code{lenet_classic} in the frozen ROCm environment do
not justify automatic transfer of numerical or runtime findings to other datasets,
architectures, or hardware. Equivalence gates establish only the registered numerical
relations of specific implementation candidates, not universal equivalence of
algorithm classes.

For QWake-FP, two boundaries are most important. First, the bounded negative C2 result
concerns a finite frozen scalar-rule family and is not proof that no zero-danger
recognizer can exist. Second, the full decision-cost surface from C1 is not the
measured marginal overhead of a particular production recognizer. These limitations
must not be repaired by post-hoc re-accounting inside the completed campaign; they
define properly separated future research directions.
